\begin{claim}{D.63}{Proved}{The exact homology characters of a negative torus centralizer}
Let $\Omega$ be the full once-punctured-torus peripheral
outer group, and let $b\in\Omega$ be a negative involution.
Keep the actual extension and peripheral labels of D.59--D.62:
\[
 \begin{gathered}
 F=\widehat F(x,y),\quad c=[x,y],\quad
 E=F\rtimes\langle v\rangle,\quad
 \alpha=\operatorname{Inn}(v)|_F,\quad
 v^2=1,\quad \alpha(c)=c^{-1},\\
 M=F_{\rm ab},\quad A=\alpha_*,\quad
 C=C_F(v),\quad gvg^{-1}=cv,\quad
 M^+=\ker(A-I)=\widehat\Z[g].
 \end{gathered}                                               \tag{D.63a}
\]
In particular, D.59 says that $[g]$ is primitive and
$C\subseteq\overline{F'}$.
Put $Z_b=C_\Omega(b)$.
There are actual continuous characters
\[
 \sigma:Z_b\longrightarrow\{1,-1\},\qquad
 \lambda:Z_b\longrightarrow\widehat\Z^\times
                                                               \tag{D.63b}
\]
such that $\lambda$ is the peripheral multiplier,
$\beta_*[g]=\sigma(\beta)[g]$, and the full homology
matrix of every $\beta\in Z_b$ is exactly
\[
 \rho(\beta)=
 \frac{\sigma(\beta)}2
 \bigl((1+\lambda(\beta))I+(1-\lambda(\beta))A\bigr)
       \quad\text{in }\operatorname{GL}_2(\widehat\Z).
                                                               \tag{D.63c}
\]
Both half-coefficients in this formula belong to
$\widehat\Z$: a full unit is odd at two, and two
is invertible at every odd prime.
In particular,
\[
          \lambda(\beta)=1
       \quad\Longrightarrow\quad
          \rho(\beta)=\sigma(\beta)I.                          \tag{D.63d}
\]

Write
\[
 \begin{split}
 \Xi_b&=\{(\sigma(\beta),\lambda(\beta)):\beta\in Z_b\},\\
 I_b&=Z_b\cap\ker\bigl(\Omega\longrightarrow
                          \operatorname{GL}_2(\widehat\Z)\bigr).
 \end{split}
\]
The exact character sequence is
\[
 1\longrightarrow I_b\longrightarrow Z_b
       \xrightarrow{\,(\sigma,\lambda)\,}\Xi_b
       \longrightarrow1.                                    \tag{D.63e}
\]
Its kernel has the following description by actual
automorphisms, with inverse maps given in the proof:
\[
 I_b\ \cong\
 J_b:=
 \{\psi\in\operatorname{Aut}(F):
          \psi(c)=c,\ \psi\alpha=\alpha\psi\}.                 \tag{D.63f}
\]
Every member of $J_b$ automatically acts trivially on $M$.
The group $I_b$ is torsion-free.
This description does not assert that $I_b$ is trivial,
or that $\Xi_b$ is the whole product in (D.63b).

Every finite subgroup $K\leq Z_b$ has the exact
faithful homology embedding
\[
 K\ \lhook\joinrel\longrightarrow\
       \{I,-I,A_2,-A_2\}\cong C_2\times C_2,
       \qquad \beta\longmapsto \rho_2(\beta).                 \tag{D.63g}
\]
Here $A_2$ is the original dyadic involution matrix.
Thus its possible abstract types are $1,C_2,C_2^2$;
a finite subgroup containing $b$ has order two or four.
This bound makes no assertion that a four-element
subgroup occurs for every $b$.
\end{claim}

\begin{proof}
\emph{The extension map for a centralizing outer class.}
The map
\[
 E\longrightarrow\operatorname{Aut}(F),\qquad
       f v^\epsilon\longmapsto
                  \operatorname{Inn}(f)\alpha^\epsilon
                                                               \tag{D.63h}
\]
identifies $E$ with the full inverse image of
$\langle b\rangle$ under
$\operatorname{Aut}(F)\to\operatorname{Out}(F)$.
Indeed an element of that inverse image is
$\operatorname{Inn}(f)\alpha^\epsilon$ by its outer label.
The map is injective because $F$ is centre-free and
$b$ has order two. It is a homeomorphism onto its image
by compactness. This also gives $C_E(F)=1$.

Let $\beta\in Z_b$, and choose an actual representative
$\psi\in\operatorname{Aut}(F)$.
Conjugation by $\psi$ preserves the inverse image in
(D.63h), since its outer class centralizes $b$.
Transport through (D.63h) gives an actual automorphism
$\Psi:E\to E$, with
$\Psi|_F=\psi$ and the identity map on $E/F=C_2$.
This constructs the required extension directly from
the centralizer; no conjugator from a standard
extension has been presumed.

Write the peripheral equation as
$\psi(c)=u c^\mu u^{-1}$, with
$\mu=\lambda(\beta)\in\widehat\Z^\times$.
Conjugate $\Psi$ by $u^{-1}$ to obtain
$\Psi(c)=c^\mu$.
The element $t=\Psi(v)$ has order two, has outside
label, and inverts $c^\mu$.
Taking its $\mu^{-1}$ power shows that it inverts $c$.
The full cusp normalizer from D.58 is
$\langle c\rangle\rtimes\langle v\rangle$.
Consequently
\[
                      t=c^zv
                 \quad(z\in\widehat\Z)
\]
with a unique exponent.
Let $i\in\{0,1\}$ be the reduction of $z$ modulo two.
Conjugating the entire extension map by
$c^{(i-z)/2}$ now gives
\[
                \Psi(c)=c^\mu,\qquad
                \Psi(v)=v_i:=c^iv.                           \tag{D.63i}
\]
The half-exponent is defined at every prime, with
even numerator at two.
All these conjugations lie in $F$, so the outer
class and its homology action have been preserved.

For completeness, (D.63i) determines the actual
representative uniquely once its outer class is fixed.
An inner comparison preserving $c^\mu$ belongs to
$C_F(c^\mu)=\langle c\rangle$, by the full peripheral
centralizer calculation. Write it as $c^a$.
It changes the exponent of $v_i$ by $2a$.
Two exponents in $\{0,1\}$ can therefore agree only
with the same $i$, and then $2a=0$ gives $a=0$.
The multiplier itself is unique by the same full
peripheral normalizer calculation.

\emph{The actual transporter determines the fixed-line sign.}
Apply $\Psi$ to $gvg^{-1}=cv$.
The resulting equation is
\[
             \Psi(g)v_i\Psi(g)^{-1}=c^\mu v_i.
                                                               \tag{D.63j}
\]
D.62 gives the full solution cosets, retaining the
specified generator $c$:
\[
 \begin{array}{c|c|c}
 i&\Psi(g)\text{ belongs to}&[\Psi(g)]\text{ in }M\\ \hline
 0&c^{(\mu-1)/2}g\,C&[g]\\
 1&c^{(\mu+1)/2}g^{-1}\,gCg^{-1}&-[g].
 \end{array}                                                  \tag{D.63k}
\]
The last column uses $C\subseteq\overline{F'}$.
Thus the action on the fixed line is specifically
$\sigma(\beta)=(-1)^i$.
Because $[g]$ generates the full direct summand
$M^+$, its actual action defines a continuous
character; composition multiplies these two signs.

\emph{The full matrix and the exact kernel.}
The peripheral multiplier is also
$\det\rho(\beta)$.
One can verify the equality at each prime $p$ by
mapping $F$ to the group of upper unitriangular
three-by-three matrices over $\Z_p$, with $x,y$
the two elementary generators.
The commutator coordinate of their images under
$\psi$ is $\det\rho_p(\beta)$.
The equation $\psi(c)=c^\mu$ makes it $\mu_p$;
a peripheral conjugator would have the same effect
because that commutator coordinate is central.
This proves the equality in every $\Z_p$ and hence
in $\widehat\Z$.

The matrices $\rho_p(\beta)$ and $A_p$ commute,
since their outer classes commute.
On $M\otimes_{\widehat\Z}\Q_p$, the matrix $A_p$
has the two distinct eigenvalues $1,-1$.
Each eigenspace has dimension one, since its
determinant is the negative multiplier $-1$.
The fixed-space eigenvalue of $\rho_p(\beta)$
is $\sigma$, and its determinant is $\mu_p$.
Its other eigenvalue is consequently $\sigma\mu_p$.
This proves (D.63c) over $\Q_p$.
The two half-coefficients are integral as explained
in the statement, and $\Z_p^2\to\Q_p^2$ is injective.
The same matrix equality therefore holds over
$\Z_p$ for every $p$, proving the displayed full
formula and (D.63d).

More explicitly, extend $[g]$ to a basis of $M$.
In that basis there is an $a\in\widehat\Z$ with
\[
 A=\begin{pmatrix}1&a\\0&-1\end{pmatrix},
 \qquad
 \rho(\beta)=
 \sigma\begin{pmatrix}
           1&a(1-\mu)/2\\0&\mu
         \end{pmatrix}.                                    \tag{D.63l}
\]
This is the same formula on the actual integral
lattice, including its possibly nonsplit dyadic
eigenspaces. It does not replace that lattice by
a direct sum of the two integral eigenlattices.

The formula proves that the kernel of
$(\sigma,\lambda)$ is exactly $I_b$.
Its image is the stated compact group $\Xi_b$,
so (D.63e) is an exact sequence of profinite groups.
In particular all commutators in $Z_b$ lie in $I_b$.
On $Z_b^1=\ker\lambda$ it restricts to the exact
sequence
\[
 1\longrightarrow I_b\longrightarrow Z_b^1
       \xrightarrow{\,\sigma\,}\sigma(Z_b^1)
       \longrightarrow1,\qquad
                 \sigma(Z_b^1)\leq\{1,-1\}.                 \tag{D.63m}
\]
No larger image is inferred.

For $\beta\in I_b$, the unique normalized
representative (D.63i) has $\mu=1$ and $i=0$.
Its extension fixes $v$ and $c$, so its restriction
belongs to $J_b$.
Conversely, a member $\psi\in J_b$ extends by
$v\mapsto v$ to $E$, and its outer class centralizes
$b$. It is already normalized with $\mu=1,i=0$.
Equations (D.63c) and (D.63k) give $\psi_*=I$,
so that outer class belongs to $I_b$.
These two maps are inverse.
Composition of two representatives fixing both
$c$ and $v$ again fixes both, so they are group
homomorphisms. They are homeomorphisms because
$J_b$ is a closed subgroup of the profinite group
$\operatorname{Aut}(F)$ and the outer projection is
a continuous bijection onto $I_b$.
Equivalently, an inner automorphism in $J_b$ would
be conjugation by $c^a$ and would require
$c^a=\alpha(c^a)=c^{-a}$, giving $a=0$.
This also verifies directly that no inner ambiguity
remains in (D.63f).

D.55 proves faithfulness of every primary homology
representation on every finite subgroup of $\Omega$.
A nonidentity torsion element of $I_b$ would generate
a finite subgroup killed by that representation.
This proves that $I_b$ is torsion-free.
It does not remove infinite-order elements of the
actual group (D.63f).

\emph{Finite centralizers.}
For $\beta$ in a finite $K\leq Z_b$, its dyadic
multiplier is a finite-order element of
$\Z_2^\times$, hence is $1$ or $-1$.
Indeed $\Z_2^\times=\{1,-1\}\times(1+4\Z_2)$,
and the latter factor is torsion-free by the
injective $2$-adic logarithm.
Formula (D.63c) at two therefore puts its matrix
in $\{I,-I,A_2,-A_2\}$.
These four matrices are distinct and form a
Klein four group: $A_2^2=I$ and
$\det A_2=-1$, which excludes $A_2=\pm I$.
D.55 makes the displayed map injective on $K$.
This proves (D.63g) and all the stated finite bounds.
The proof leaves exactly the kernel in (D.63f)
and the actual image $\Xi_b$ to be determined.
\end{proof}
